\subsection{Matched deployment: k3 gains survive fresh contexts}
\label{sec:matchedbridge}
A new matched study separates the effects of speculative depth and context lifecycle within one runtime~\cite{rtxnativebridge}. It uses the original twelve exposed Qwen3 inputs, unchanged model/tokenizer bytes, context1024, target UB4, draft UB1 and cap128/native EOG. One executable supplies all four $k=1,3$ by resident/fresh cells and the PQ\_A/PQ\_B/QB arms. All 204 qualification requests preserve complete outputs and speculative state; corresponding PQ/QB decision vectors agree, and every output also matches the historical E2 reference.

Six independently loaded process sessions complete 864 measurements and 72 warmups. The order balances cells and arms within input/session blocks. Both models remain loaded. Resident requests clear their existing contexts; fresh requests create and destroy contexts inside the service timer. The shared timer includes prefill, generation, sampling, synchronization and result construction. Heavy capture and additional backend counters are disabled.

\begin{table*}[t]\centering\small
\caption{Matched Qwen3 deployment bridge. Ratios compare QB with the arithmetic mean of identical PQ controls; aggregation gives equal weight to inputs and process sessions. Intervals are the registered Bonferroni-adjusted session-bootstrap intervals for the four-cell family. The practical screen also requires reduction above the A/A threshold and improvement in every session.}
\label{tab:nativebridge}
\begin{tabular}{@{}rlll@{}}\toprule
$k$ & Lifecycle & QB/dual-PQ [simultaneous interval] & Practical gain\\\midrule
1 & Resident & 0.986617 [0.965043, 1.006956] & Not established\\
3 & Resident & 0.937404 [0.922603, 0.949354] & Established\\
1 & Fresh & 0.973534 [0.955904, 0.996512] & Not established\\
3 & Fresh & 0.946409 [0.940517, 0.953312] & Established\\
\bottomrule\end{tabular}
\end{table*}


QB reduces service time by 6.26\% at resident $k=3$ and 5.36\% at fresh $k=3$ (Table~\ref{tab:nativebridge}). Both pass the registered practical screen; neither $k=1$ cell passes. The $k=3$ point reductions exceed their respective 4.125\% and 4.388\% input-aggregate A/A imbalances and favor QB in every session. The four primary intervals resample whole process sessions and use a Bonferroni adjustment; the twelve inputs remain a fixed panel.

The layout benefit therefore survives fresh-context service at $k=3$. Secondary paired contrasts favor a stronger layout effect at $k=3$ in both lifecycles, while each individual lifecycle contrast includes parity. Changing $k$ changes the speculative transaction as well as query width. These comparisons identify a deployment-dependent layout effect, not a pure width intervention or a CUDA Graph mechanism. The new driver also does not reconstruct every difference between the older HTTP E2 and fresh-service contracts.

\subsection{Transfer to previously unexposed inputs}
\label{sec:transfer}
A separate study retains the three-family $k=3$ resident HTTP contract and tests twelve previously unexposed natural inputs per family~\cite{rtxtransferrecords}. The frozen panel covers formatted prompt-length strata 64--192, 256--448 and 640--832 within context1024. Short inputs are GSM8K mathematics; longer inputs are CodeContests code. Length and domain therefore vary together.

\begin{table*}[t]\centering\small
\caption{Transfer to previously unexposed inputs within context1024. Each family has a registered 12-input panel. Qwen3 and SmolLM2 complete all inputs; the Llama reference failure stops that family before cost measurement. The layout screen uses a 2\% floor, input-aggregate A/A and every-session direction; no population interval is assigned.}
\label{tab:heldouttransfer}
\begin{tabular}{@{}llll@{}}\toprule
Family & Reference qualification & QB/dual-PQ (screen) & QB/R\\\midrule
Qwen3 & 12/12 inputs pass & 0.967781 (not established) & 0.560577\\
SmolLM2 & 12/12 inputs pass & 0.930056 (pass) & 1.007007\\
Llama & PQ and QB both depart from R & Not run & Not run\\
\bottomrule\end{tabular}
\end{table*}


Qwen3 and SmolLM2 pass all twelve inputs. Four sessions per family supply 384 measurements and 384 warmups in total. SmolLM2's 6.99\% QB layout reduction passes the finite-panel practical screen, yet QB/R is 1.007007: the complete speculative request remains slightly slower than its serial reference. Qwen3's 3.22\% reduction is smaller than its 4.78\% A/A threshold. Layout value and profitable speculation remain separate questions.

Llama provides a reference-transfer counterexample. On CodeContests/valid/108, with 737 formatted prompt tokens, PQ and QB agree with each other in outputs, full decision vectors and causal trace, but both depart from deployed R at zero-based output index81. The remaining three Llama inputs and all Llama timing slots stay unexecuted. PQ/QB correspondence therefore does not establish reference preservation on a new input. The tested native-fast comparator also fails reference qualification in all three families; no qualified timing victory over that route is claimed. \sref{supp:followup} gives accounting and source access.
